\chapter{The Virtual Power Grid}
\label{ch:grid}

Running catenary wire across a city is honest work and is not always practical. The
\textbf{virtual power grid} moves Immersive Flux between places that are not connected by any wire
at all --- across a mountain, across a chunk border, across dimensions.

\begin{iefork}
Not in stock Immersive Engineering. It takes the idea Flux Networks made popular and rebuilds it in
IE's idiom: sheet-metal service boxes, a console you walk to, and named segments rather than
per-player networks.
\end{iefork}

\section{Feed Units and Service Units}

Power still has to enter and leave somewhere. It enters at a \textbf{Grid Feed Unit} and leaves at a
\textbf{Grid Service Unit}. Both bolt to any solid surface, including engineering posts.

These boxes exchange power with the block they are bolted to, exactly as a machine does. They do
\emph{not} accept wires directly --- put a connector against one if you want it to feed a wire
network. A Feed Unit takes power out of the world; a Service Unit puts it back. Both carry a
terminal post on the front, which is where wiring goes.

\begin{ietip}
A wire connector \emph{beside} a Service Unit is fed as well as one bolted to it, whichever way the
connector faces. IE connectors normally take power on one side only --- the block they are mounted
on --- and a connector on the wall next to a unit would otherwise sit there touching live hardware
and doing nothing.
\end{ietip}

\section{Segments}

The grid is not one pool. It is divided into \textbf{segments}: named groups of boxes, each with its
own switch, its own transfer limits and its own statistics. A segment is the thing you switch off
when you want the east side of town dark; the boxes on it are the sockets.

A box that has not been assigned to a segment sits with an unlit lamp doing nothing. Assigning it is
the whole configuration step.

Each segment carries a \textbf{priority} order and a \textbf{transfer cap} per device, so you decide
what is served first when there is not enough to go round. Anything marked \textbf{critical} is
served before everything else --- the pumps, the blast furnace, the hospital.

\section{The Grid Management Console}

A $2\times2$ wall of four different blocks, struck anywhere on its front face with an Engineer's
Hammer:

\begin{center}
\small
\begin{tabular}{@{}lll@{}}
\toprule
 & Left & Right \\
\midrule
Upper & Console Housing (the terminal) & Redstone Engineering Block \\
Lower & Light Engineering Block & Heavy Engineering Block \\
\bottomrule
\end{tabular}
\end{center}

\noindent The terminal is the screen, the redstone block is the instrument rack beside it, and the
light and heavy blocks are the desk and the power cabinet under them. The finished console is drawn
as one piece --- one desk, one monitor spanning both blocks --- and breaking it returns the four
components.

Six tabs: an overview of every segment, the segment editor, the device list, the failover editor,
live statistics, and settings. Right-clicking any box also opens its own small panel, so a freshly
placed unit can be assigned without walking back.

\section{Failover}

A segment may name another as its \textbf{backup}. If the first cannot meet its load, the second
covers it. With \emph{top-up} enabled a backup also covers ordinary shortfalls, not only outages;
with it off, backups engage on outage or trip alone.

\section{Breakers and schedules}

If breakers are enabled, a segment held at its output ceiling long enough will \textbf{trip} and
latch off until somebody flips it back on. That is a real switch on the Segments tab.

A segment may also carry a \textbf{schedule}: two times of day between which it may run. The
schedule is a gate, never a second switch --- it can hold a segment down but never raise one a
player switched off, because otherwise the console toggle and the clock fight every dusk and
whichever ran last wins.

\section{The Grid Signal Unit}

Bridges a segment to redstone in either direction. In \emph{input} mode a signal holds its segment
off, which makes any lever, daylight sensor or logic circuit an external kill switch; inverted, it
demands a keep-alive signal instead, which is a dead-man's switch. In \emph{output} mode it reports
whether its segment is up --- or, inverted, whether it is down, which is a fault lamp.

\section{The Engineer's Network Terminal}

A held item showing every segment and every fluid main, from anywhere, and changing neither.

\begin{ietip}
The terminal is read-only by construction rather than by convention --- there is no path from it to
a change, whatever the client sends. Changes are made at a console.
\end{ietip}

\section{Chunk loading}

A device may be marked to keep its chunk loaded, within a server-wide budget. Fittings sharing a
chunk cost one between them. A disabled device stops holding its chunk.

\section{City mode}

See \autoref{ch:citymode}. In short: segments stop accounting for flux and switch to presence. A
segment is energised or it is not, and its Service Units deliver freely.
